% Pursuing Stacks (cotes 134-2 to 134-8): archive cote 134-7, site page 98 (PDF page 99), typescript page 540.
% typescript. Draft compared: ps6.tex:4949-5011.
% AI-transcribed by Claude Opus 5.5 (claude-opus-5-5, adaptive thinking, effort high), Message Batches API, batch ps4 (msgbatch_012Aoc1NuskATPsyBeKSYK7J), 2026-11-12; not yet checked by eye.
\scanpage{134-7/98}
\mspage{540}

\xout{$p_*(\underline{O}_X)$ def def $p_*(\underline{O}_X)$}
A is the \xout{k} k-module (a k-algebra as a matter of fact) such\xout{t}that
\margin{\uline{$\Gamma(S,p_*(\underline{O}_X))$} (circled)}
\[ \xout{p_*(\underline{O}_X)} \quad \overset{\hand{\sim}}{A} \simeq p_*(\underline{O}_X) \quad , \hand{i.e. $A=\Gamma(S,p_*(\underline{O}_X))\simeq\Gamma(X,\underline{O}_X)$} \]
the last member of (*) is
\[ \mathrm{Hom}_{\underline{O}_S}(\overset{\hand{\sim}}{M} , \overset{\hand{\sim}}{A}) \simeq \mathrm{Hom}_k(M,A) \hand{\simeq} \mathrm{Hom}_{\underline{O}_k}(V(A), V(M)) \quad , \quad \xout{xx} \]
and finally we get
\[ (14) \qquad \mathrm{Hom}(X,V(M)) \hand{\simeq} \mathrm{Hom}\xout{k}_{\underline{O}_k}(V(A),V(M)) \qquad . \]
This shows that the \xout{xxxxxxxx}\inserted[typed]{forgetful} functor from the category $\mathrm{Uab}_V$ of all
vector bundles over k to the category of all \xout{xx} k-schemes which are qua-%
si-compact and quasi-se\xout{p}parated, admits a left adjoint $L_V$, where
\margin{\uline{$\Gamma(X,\underline{O}_X)$.} (circled)}
\[ (15) \qquad L_{\hand{V}}(X) = V(A) \quad , \quad \text{where} \quad A = \Gamma(S, p_{X\,*}(\underline{O}_X)) \hand{$\simeq \Gamma(X,\underline{O}_X)$} \]
In case X is affine, A is just the affine ring of \xout{x}X , which is a k-alge-%
bra, from which we retain only (in formula (15)) the k-module structure.

Thus, as far as the notion of k-linearization goes, V-bundles
behave in a considerably nice\xout{s}r way than W-bundles, without any need to
go over to proob\xout{g}jects. \xout{One may wonder} Th\xout{i}us it may be preferable to
work with V-bundles rather than W-bundles. We may wonder at this point
why not \xout{take}\inserted[typed]{admit} then as``bundle\xout{s}''any \inserted[typed]{k-}scheme X which is quasicompact and
quasiseparated, or at any rate any affine k-scheme, as these can be
quite conveniently described in terms of k-algebras. Thus\hand{,} semisimplici-%
al affine schemes over k just correspond to \uline{co}-semisimplicial commutati-%
ve algebras over k, and likewise for maps - and linearization just
corresponds to forgetting the algebra structure in this co-semisimplicia
object, and retain the structure of a co-semisimplicial k-module,
which corresponds dually to a semisimplicial vector bundle. This is
a perfectly simple relationship - why bother about restricting the
``bundles'' from arbitrary affine k-schemes to those which are isomorphic
to a vector bundle ?

We should remember here our initial motivation, which was, in
case when k=$\underline{Z}$, to get a category of \xout{``homotopy} ``schematic'' homotopy
types as close as possible to the usual one. One plausible way of
achi\xout{f}eving this is by restricting the notion of a bundle the more we
can, so as to get still the possibility of ``schematization'' for a very
sizeable bunch of ordinary homotopy types. Postnikov \unsure{dévisaage|dévissage} than
suggested to work with so-called``unipotent bundles'', and it was almost

\margin{98}
